\section{Introduction}
Passwords constitute the primary authentication factor for the vast majority of currently deployed web applications. The prevalence of attacks like credential stuffing \cite{credential_stuffing} has made safe password management amongst the most critical security tasks that an average user faces today. Accordingly, a variety of popular password management tools have emerged to help users address this issue in a secure and convenient manner.

The most commonly used password management solutions are cloud-based applications like LastPass \cite{lastpass_arch}, Dashlane \cite{dashlane_arch}, and 1Password \cite{1password_arch}. These systems store user credentials in a centralized vault, typically encrypted using a key derived from the user's master password via PBKDF2 \cite{rfc2898}. In doing so, they afford users the convenience of accessing their accounts from any device, but also create a central point of failure that, if compromised, could reveal all of a user's passwords. Indeed, LastPass alone has experienced 8 major security incidents \cite{lastpass_security}, including a recent total data breach of stored credentials \cite{toubba_notice_2022}.

Alternatively, open-source password management software like Bitwarden \cite{bitwarden} and KeePassX \cite{keepassx} has emerged in part to eliminate the central point of failure created by cloud-based password managers. While popular amongst experienced users, these solutions have failed to achieve widespread adoption due their relative lack of usability. In particular, synchronizing data between devices is often a manual and cumbersome process.

\eject

For over two decades, deterministic password generators (DPGs) have been proposed as a substitute for all forms of password management \cite{karp2002site}. DPGs like PwdHash \cite{ross2005stronger} apply a cryptographic hash function to a user's master password and the domain name of a website to generate a unique pseudorandom password for each domain. As such, DPGs theoretically work seamlessly across devices without the need for synchronization, offering the security of a client-side password manager with the flexibility of a cloud-based solution. 

In practice, DPGs have been widely criticized for having a variety of security, privacy, and usability flaws \cite{noauthor_4_2016} that have thus far seriously hindered their adoption.
In this paper, we present a detailed analysis of 45 existing DPGs and summarize the main issues hindering their adoption.
In particular, current DPG schemes allow a user's master password to be directly attacked in the event that any of a user's generated passwords are compromised. Most DPGs also lack the flexibility to support password rotation or complex password policies.

Based on these findings, we present a new multi-factor deterministic password generator (MFDPG), which aims to rectify the shortcomings of existing DPGs by incorporating multi-factor key derivation \cite{mfkdf} into the password management process. In doing so, MFDPG effectively allows users to unilaterally upgrade password-only websites to support strong multi-factor authentication like TOTP \cite{rfc6238} and YubiKey \cite{yubikey}.
We further propose novel algorithmic solutions that facilitate password rotation and complex password policy compliance without leaking service usage patterns.
By presenting a truly secure and practical design, we hope to revive the notion of DPGs as a viable alternative to password management. \\

\noindent \textbf{Contributions:}
\vspace{-0.5em}
\begin{enumerate}[leftmargin=*]
    \item We analyze 45 existing DPGs to asses the security, privacy, and usability issues hindering their adoption (\S\ref{sec:survey}).
    \item We present MFDPG, a practical multi-factor authenticated DPG with zero client or server-side secret storage (\S\ref{sec:approach}).
    \begin{enumerate}
        \item MFDPG uses multi-factor key derivation to harden the system against attacks on the master password (\S\ref{sec:mfdpg}).
        \item We use Cuckoo filters to solve the revocation problem without leaking private account information (\S\ref{sec:revocation}).
        \item Our novel password generation algorithm supports any regular password policy using DFA traversal (\S\ref{sec:generation}).
    \end{enumerate}
    \item We evaluated MFDPG to verify its compatibility with all of the 100 most popular existing web applications (\S\ref{sec:evaluation}).
    \item MFDPG has the further effect of progressively upgrading any password-based website to support strong MFA (\S\ref{sec:discussion}).
\end{enumerate}

\eject